\documentclass[10pt,a4paper]{amsart}
\usepackage{amsmath,amssymb,amsthm,graphicx}
\usepackage{verbatim}
\newenvironment{reference}{\comment}{\endcomment}
%\newenvironment{reference}{}{}
\newenvironment{slogan}{\comment}{\endcomment}
\newenvironment{history}{\comment}{\endcomment}

% For commutative diagrams we use Xy-pic
\usepackage[all]{xy}

% We use 2cell for 2-commutative diagrams.
\xyoption{2cell}
% Theorem environments.
%
\theoremstyle{plain}
\newtheorem{theorem}[subsection]{Theorem}
\newtheorem{proposition}[subsection]{Proposition}
\newtheorem{lemma}[subsection]{Lemma}

\theoremstyle{definition}
\newtheorem{definition}[subsection]{Definition}
\newtheorem{example}[subsection]{Example}
\newtheorem{exercise}[subsection]{Exercise}
\newtheorem{situation}[subsection]{Situation}

\theoremstyle{remark}
\newtheorem{remark}[subsection]{Remark}
\newtheorem{remarks}[subsection]{Remarks}

\numberwithin{equation}{subsection}

% Macros
%
\def\lim{\mathop{\mathrm{lim}}\nolimits}
\def\colim{\mathop{\mathrm{colim}}\nolimits}
\def\Spec{\mathop{\mathrm{Spec}}}
\def\Hom{\mathop{\mathrm{Hom}}\nolimits}
\def\Ext{\mathop{\mathrm{Ext}}\nolimits}
\def\SheafHom{\mathop{\mathcal{H}\!\mathit{om}}\nolimits}
\def\SheafExt{\mathop{\mathcal{E}\!\mathit{xt}}\nolimits}
\def\Sch{\mathit{Sch}}
\def\Mor{\mathop{\mathrm{Mor}}\nolimits}
\def\Ob{\mathop{\mathrm{Ob}}\nolimits}
\def\Sh{\mathop{\mathit{Sh}}\nolimits}
\def\NL{\mathop{N\!L}\nolimits}
\def\CH{\mathop{\mathrm{CH}}\nolimits}
\def\proetale{{pro\text{-}\acute{e}tale}}
\def\etale{{\acute{e}tale}}
\def\QCoh{\mathit{QCoh}}
\def\Ker{\mathop{\mathrm{Ker}}}
\def\Im{\mathop{\mathrm{Im}}}
\def\Coker{\mathop{\mathrm{Coker}}}
\def\Coim{\mathop{\mathrm{Coim}}}

% Boxtimes
%
\DeclareMathSymbol{\boxtimes}{\mathbin}{AMSa}{"02}

%
% Macros for moduli stacks/spaces
%
\def\QCohstack{\mathcal{QC}\!\mathit{oh}}
\def\Cohstack{\mathcal{C}\!\mathit{oh}}
\def\Spacesstack{\mathcal{S}\!\mathit{paces}}
\def\Quotfunctor{\mathrm{Quot}}
\def\Hilbfunctor{\mathrm{Hilb}}
\def\Curvesstack{\mathcal{C}\!\mathit{urves}}
\def\Polarizedstack{\mathcal{P}\!\mathit{olarized}}
\def\Complexesstack{\mathcal{C}\!\mathit{omplexes}}
% \Pic is the operator that assigns to X its picard group, usage \Pic(X)
% \Picardstack_{X/B} denotes the Picard stack of X over B
% \Picardfunctor_{X/B} denotes the Picard functor of X over B
\def\Pic{\mathop{\mathrm{Pic}}\nolimits}
\def\Picardstack{\mathcal{P}\!\mathit{ic}}
\def\Picardfunctor{\mathrm{Pic}}
\def\Deformationcategory{\mathcal{D}\!\mathit{ef}}
\usepackage{cleveref}
\crefname{originallemma}{original lemma}{original lemmas}
\Crefname{originallemma}{Original lemma}{Original lemmas}
\makeatletter
\newlabel{more-morphisms-lemma-integral-closure-smooth-pullback}{{081K}{source}{}{sourcepointer-1}{}}
\newlabel{more-morphisms-lemma-integral-closure-smooth-pullback@cref}{{[originallemma][0][]081K}{[1][0][]source}{}{}{}}
\makeatother
\crefname{originallemma}{original lemma}{original lemmas}
\Crefname{originallemma}{Original lemma}{Original lemmas}
\makeatletter
\newlabel{morphisms-lemma-integral-closure}{{035F}{source}{}{sourcepointer-2}{}}
\newlabel{morphisms-lemma-integral-closure@cref}{{[originallemma][0][]035F}{[1][0][]source}{}{}{}}
\makeatother
\crefname{originalsection}{original section}{original sections}
\Crefname{originalsection}{Original section}{Original sections}
\makeatletter
\newlabel{spaces-properties-section-quasi-coherent}{{03G5}{source}{}{sourcepointer-3}{}}
\newlabel{spaces-properties-section-quasi-coherent@cref}{{[originalsection][0][]03G5}{[1][0][]source}{}{}{}}
\makeatother
\title{Existence of the integral closure of the structure sheaf inside a quasi-coherent algebra on an algebraic space}
\author{The Stacks Project Authors}
\date{stacks-project:spaces-morphisms.tex:local-snapshot-20261012; GNU FDL 1.2 or later}
\begin{document}
\maketitle
\setcounter{subsection}{1}
\noindent\textit{Editorial note.} All delivered mathematical text is the author's own, copied from the pinned Stacks Project chapter file spaces-morphisms.tex; no statement or proof step was rewritten, reordered or removed.\par\medskip
\noindent\textit{Editorial note.} The proof is a gluing argument that reduces to the scheme case on each object of the etale site; the compatibility isomorphism is what makes the construction canonical.\par\medskip
\noindent\textit{Editorial note.} The three references of the counted spans are recorded as non-load-bearing navigation pointers; the argument that consumes them stays inside the counted proof span.\par\medskip
\noindent\textit{Editorial note.} Preamble: the three native excerpt spans of sources/stacks-spaces-morphisms/preamble.tex (lines 8-13, 14-18 and 154-217). The chapter file itself carries no preamble text.\par\medskip
\noindent\textit{Allowed standard prerequisites (editorial).}\begin{itemize}
\item The integral closure of the structure sheaf inside a quasi-coherent algebra on a scheme. Morphisms, Lemma morphisms-lemma-integral-closure, used at proof lines 9946-9952. Is the scheme case applied on each U in X\_etale.
\item Integral closures commute with smooth (hence etale) pullback. More on Morphisms, Lemma more-morphisms-lemma-integral-closure-smooth-pullback, used at proof lines 9958-9962. Supplies the canonical isomorphism making the local data glue.
\item Quasi-coherent modules on X\_etale described by their values on affine etale objects. Properties of Spaces, Section spaces-properties-section-quasi-coherent, used at proof lines 9955-9957 and 9967-9972. Turns the local quasi-coherence into quasi-coherence on X.
\item Restriction of a quasi-coherent sheaf along a morphism of X\_etale. used at proof lines 9952-9958. Identifies A|\_\{U'\} with the pullback of A|\_U.
\end{itemize}\medskip
\ifdefined\hypertarget\hypertarget{sourcepointer-1}{}\fi
\noindent\textit{Source reference.} Original lemma 081K: focused-seeds/AGP/sources/stacks-snapshot-20261011/tags.txt, source lines 10220--10220. Navigation pointer into the Stacks Project; the standard fact it supplies is applied inside the counted proof span and the argument resting on it stays there.\par\medskip
\ifdefined\hypertarget\hypertarget{sourcepointer-2}{}\fi
\noindent\textit{Source reference.} Original lemma 035F: focused-seeds/AGP/sources/stacks-snapshot-20261011/tags.txt, source lines 4126--4126. Navigation pointer into the Stacks Project; the standard fact it supplies is applied inside the counted proof span and the argument resting on it stays there.\par\medskip
\ifdefined\hypertarget\hypertarget{sourcepointer-3}{}\fi
\noindent\textit{Source reference.} Original section 03G5: focused-seeds/AGP/sources/stacks-snapshot-20261011/tags.txt, source lines 4516--4516. Navigation pointer into the Stacks Project; the standard fact it supplies is applied inside the counted proof span and the argument resting on it stays there.\par\medskip
\begin{lemma}
\label{lemma-integral-closure}
Let $S$ be a scheme. Let $X$ be an algebraic space over $S$.
Let $\mathcal{A}$ be a quasi-coherent sheaf of $\mathcal{O}_X$-algebras.
There exists a quasi-coherent sheaf of $\mathcal{O}_X$-algebras
$\mathcal{A}' \subset \mathcal{A}$ such that
for any affine object $U$ of $X_\etale$ the ring
$\mathcal{A}'(U) \subset \mathcal{A}(U)$ is
the integral closure of $\mathcal{O}_X(U)$ in $\mathcal{A}(U)$.
\end{lemma}

\begin{proof}
Let $U$ be an object of $X_\etale$. Then $U$ is a scheme. Denote
$\mathcal{A}|_U$ the restriction to the Zariski site. Then
$\mathcal{A}|_U$ is a quasi-coherent sheaf of $\mathcal{O}_U$-algebras
hence we can apply Morphisms, Lemma \ref{morphisms-lemma-integral-closure}
to find a quasi-coherent subalgebra $\mathcal{A}'_U \subset \mathcal{A}|_U$
such that the value of $\mathcal{A}'_U$ on any affine open $W \subset U$
is as given in the statement of the lemma. If $f : U' \to U$ is a morphism
in $X_\etale$, then $\mathcal{A}|_{U'} = f^*(\mathcal{A}|_U)$
where $f^*$ means pullback by the morphism $f$ in the Zariski topology;
this holds because $\mathcal{A}$ is quasi-coherent (see introduction
to Properties of Spaces, Section \ref{spaces-properties-section-quasi-coherent}
and the references to the discussion in the chapter on
descent on schemes). Since $f$ is \'etale we find that
More on Morphisms, Lemma
\ref{more-morphisms-lemma-integral-closure-smooth-pullback}
says that we get a canonical isomorphism
$f^*(\mathcal{A}'_U) = \mathcal{A}'_{U'}$.
This immediately tells us that we obtain a sub presheaf
$\mathcal{A}' \subset \mathcal{A}$ of $\mathcal{O}_X$-algebras
over $X_\etale$ which is a sheaf for the Zariski
topology and has the right values on affine objects.
But the fact that each $\mathcal{A}'_U$ is quasi-coherent
on the scheme $U$ and that for $f : U' \to U$ \'etale we have
$\mathcal{A}'_{U'} = f^*(\mathcal{A}'_U)$ implies that $\mathcal{A}'$
is quasi-coherent on $X_\etale$ as well (as this is a local property
and we have the references above describing quasi-coherent modules
on $U_\etale$ in exactly this manner).
\end{proof}

\end{document}
